\subsection{Violencia letal: mejora general y diferencias por género}
La serie de Infosegura registra una tasa total de homicidios de 106,3
por cien mil habitantes en 2015 y de 20,9 en 2020. La proyección rotulada para
2021 es 20,8. En mujeres, la tasa correspondiente es 17,3 en 2015, 3,8 en
2020 y 4,9 en la proyección de 2021. La serie de violencia de género informa tasas de
feminicidio de 2,2 en 2020 y 3,1 en su proyección de 2021. Se preserva la
condición proyectada del último punto, sin presentarlo como cierre anual.

\begin{figure}[htbp]
\centering
\begin{minipage}[t]{0.49\textwidth}\vspace{0pt}\centering
\begin{tikzpicture}
\begin{axis}[width=0.9\linewidth,height=5.7cm,title={Homicidios: total},
 xmin=2013,xmax=2021,ymin=0,ymax=120,xtick={2013,2015,2017,2019,2021},
 xticklabel style={rotate=45,anchor=east},scaled x ticks=false,
 /pgf/number format/1000 sep={},ylabel={Por cien mil habitantes},
 tick label style={font=\scriptsize},title style={font=\small},
 grid=major,grid style={gray!15},filter discard warning=false]
\addplot[petroleo,thick,mark=*,restrict x to domain=2013:2020]
 table[x=anio,y=total,col sep=comma]{data/infosegura_homicidios.csv};
\addplot[petroleo,thick,dashed,mark=o,restrict x to domain=2020:2021]
 table[x=anio,y=total,col sep=comma]{data/infosegura_homicidios.csv};
\end{axis}\end{tikzpicture}
\end{minipage}\hfill
\begin{minipage}[t]{0.49\textwidth}\vspace{0pt}\centering
\begin{tikzpicture}
\begin{axis}[width=0.9\linewidth,height=5.7cm,title={Muertes de mujeres},
 xmin=2013,xmax=2021,ymin=0,ymax=20,xtick={2013,2015,2017,2019,2021},
 xticklabel style={rotate=45,anchor=east},scaled x ticks=false,
 /pgf/number format/1000 sep={},ylabel={Por cien mil mujeres},
 tick label style={font=\scriptsize},title style={font=\small},
 grid=major,grid style={gray!15},filter discard warning=false,
 legend style={font=\scriptsize,at={(0.5,-0.45)},anchor=north,draw=none}]
\addplot[petroleo,thick,mark=*,restrict x to domain=2013:2020]
 table[x=anio,y=mujeres,col sep=comma]{data/infosegura_homicidios.csv};
\addlegendentry{Homicidios de mujeres}
\addplot[ocre,thick,mark=square*,restrict x to domain=2013:2020]
 table[x=anio,y=feminicidios,col sep=comma]{data/infosegura_homicidios.csv};
\addlegendentry{Feminicidios}
\addplot[petroleo,thick,dashed,mark=o,forget plot,restrict x to domain=2020:2021]
 table[x=anio,y=mujeres,col sep=comma]{data/infosegura_homicidios.csv};
\addplot[ocre,thick,dashed,mark=square,forget plot,restrict x to domain=2020:2021]
 table[x=anio,y=feminicidios,col sep=comma]{data/infosegura_homicidios.csv};
\end{axis}\end{tikzpicture}
\end{minipage}
\caption{Violencia letal en Infosegura: tasas y proyecciones, 2013-2021}\label{fig:infosegura_letal}
\fuente{PNUD/Infosegura (2021), con datos de mesas técnicas y DIGESTYC. El tramo discontinuo termina en la \textbf{proyección de 2021}. Los paneles usan escalas y poblaciones distintas. El total de homicidios de mujeres incluye feminicidios; las dos series del panel derecho no se suman.}
\end{figure}

El contraste es relevante para el argumento del artículo: una mejora de la
tasa nacional no permite asumir una evolución homogénea por género. MICS
agrega información sobre prácticas y normas que los registros letales no
describen. La coexistencia documental de disciplina violenta en 2014 y
reducción de homicidios después de 2015 no demuestra continuidad individual,
persistencia cuantificada de violencia familiar en 2021 ni efectos de una
política específica.

El cuadro \ref{tab:cambio_letal} cuantifica el descenso entre dos años
observados. La tasa total cae 80,3\,\%, pero la de feminicidios lo hace
72,3\,\%, aproximadamente ocho puntos menos en términos de reducción
relativa. Se trata de porcentajes de cambio de tasas, distintos de puntos
porcentuales de prevalencia. Las bases iniciales son diferentes; una menor
caída relativa no implica mayor número absoluto de hechos. La clasificación
de feminicidio exige además considerar prácticas de registro y tipificación.

\begin{table}[htbp]\centering\small
\caption{Cambio observado de tasas de violencia letal, 2015--2020}\label{tab:cambio_letal}
\begin{tabularx}{\textwidth}{@{}Yrrrr@{}}\toprule
Serie & 2015 & 2020 & Diferencia & Cambio (\%)\\\midrule
Homicidios, total & 106,3 & 20,9 & -85,4 & -80,3\\
Homicidios, hombres & 206,7 & 39,7 & -167,0 & -80,8\\
Homicidios, mujeres & 17,3 & 3,8 & -13,4 & -77,8\\
Feminicidios & 7,8 & 2,2 & -5,7 & -72,3\\
\bottomrule

\end{tabularx}
\fuente{Cálculos del artículo sobre las series de PNUD/Infosegura (2021), con precisión disponible antes del redondeo a una décima. Tasas por cien mil habitantes del grupo correspondiente. Diferencia en unidades de tasa; cambio: $100(r_{2020}/r_{2015}-1)$. No interviene la proyección de 2021.}
\end{table}

Los agregados de Infosegura también incluyen promedios mensuales de denuncias
intrafamiliares de 96,1, 77,9 y 105,4 para 2019, 2020 y 2021, respectivamente.
La fuente no precisa el corte de 2021. Se conservan como promedios reportados,
sin anualizarlos ni calcular con ellos una variación de
incidencia familiar. La ecuación \ref{eq:registro} ayuda a explicar por qué
una denuncia no equivale a una medición exhaustiva de violencia ocurrida.

La nota institucional de octubre de 2021 mantiene una proyección nacional
de 20,8 homicidios por cien mil habitantes y advierte una reducción desigual
por sexo \citep{pnud2021octubre}. Esa lectura respalda desagregar las
mejoras; no convierte una proyección en cierre anual ni prueba un mecanismo
familiar. Sus cifras resumidas de denuncias no se interpretan como totales
anuales comparables: aquí se conservan las unidades de promedio mensual.

\subsection{Victimización por delincuencia y exposición territorial}
La serie de victimización en los últimos doce meses registra 14,0\,\%
en 2017, 13,9\,\% en 2018, 10,8\,\% en 2019 y 10,4\,\% en 2021. En este
último año los valores son 10,7\,\% para hombres y 10,2\,\% para mujeres,
mientras área urbana y rural registran 12,7\,\% y 6,8\,\%. Se conservan las
definiciones de la fuente: estas cifras de delincuencia no sustituyen
la prevalencia de disciplina o violencia de pareja.

\begin{figure}[htbp]
\centering
\begin{minipage}[t]{0.49\textwidth}\vspace{0pt}\centering
\begin{tikzpicture}
\begin{axis}[width=0.9\linewidth,height=5.4cm,title={Por sexo},
 xmin=2017,xmax=2021,ymin=0,ymax=20,xtick={2017,2018,2019,2021},
 scaled x ticks=false,/pgf/number format/1000 sep={},ylabel={Porcentaje},
 grid=major,grid style={gray!15},tick label style={font=\scriptsize},
 title style={font=\small},legend style={font=\scriptsize,at={(0.5,-0.25)},anchor=north,draw=none}]
\addplot[petroleo,thick,mark=*] table[x=anio,y=hombres,col sep=comma]{data/infosegura_victimizacion.csv};
\addlegendentry{Hombres}
\addplot[ocre,thick,mark=square*] table[x=anio,y=mujeres,col sep=comma]{data/infosegura_victimizacion.csv};
\addlegendentry{Mujeres}
\end{axis}\end{tikzpicture}
\end{minipage}\hfill
\begin{minipage}[t]{0.49\textwidth}\vspace{0pt}\centering
\begin{tikzpicture}
\begin{axis}[width=0.9\linewidth,height=5.4cm,title={Por residencia},
 xmin=2017,xmax=2021,ymin=0,ymax=20,xtick={2017,2018,2019,2021},
 scaled x ticks=false,/pgf/number format/1000 sep={},ylabel={Porcentaje},
 grid=major,grid style={gray!15},tick label style={font=\scriptsize},
 title style={font=\small},legend style={font=\scriptsize,at={(0.5,-0.25)},anchor=north,draw=none}]
\addplot[petroleo,thick,mark=*] table[x=anio,y=urbana,col sep=comma]{data/infosegura_victimizacion.csv};
\addlegendentry{Urbana}
\addplot[ocre,thick,mark=square*] table[x=anio,y=rural,col sep=comma]{data/infosegura_victimizacion.csv};
\addlegendentry{Rural}
\end{axis}\end{tikzpicture}
\end{minipage}
\caption{Victimización por delincuencia en los últimos doce meses}\label{fig:infosegura_vict}
\fuente{PNUD/Infosegura (2021), con encuestas de Cultura de Paz. Años observados: 2017, 2018, 2019 y 2021; no se imputa 2020. Los agregados no explicitan la edad del universo ni incluyen IC; las diferencias no se presentan como significativas.}
\end{figure}

La brecha entre hombres y mujeres pasa de 5,39 puntos en 2018 a 1,39 en
2019 y 0,49 en 2021. La urbana-rural, en cambio, pasa de 2,43 en 2019 a
5,86 en 2021; esta última distancia se aproxima a los 6,37 puntos de 2018.
Por tanto, la caída del total no describe por sí sola la distribución
territorial. Al no disponer de errores estándar y covarianzas entre
dominios, la reducción de una brecha y ampliación de otra son contrastes
puntuales, no evidencia de cambios estadísticamente significativos. Un
diagnóstico de cobertura debe conservar ambos ejes de desagregación.

La brecha urbana-rural de victimización tiene el mismo signo que la diferencia
descriptiva de disciplina de MICS y distinto signo que la justificación de
violencia contra la esposa. Esa observación no es una correlación estimada:
son años, unidades y experiencias diferentes. Su utilidad consiste en evitar
que un indicador territorial se convierta en sustituto automático de todos
los demás y en orientar una agenda de desagregación comparable.

\subsection{Condiciones materiales, prioridades y percepción}
Los indicadores de opinión sitúan economía y seguridad como principal problema del país
en 46,9\,\% y 26,6\,\% en 2021. En 2019 informa 32,3\,\% y 59,8\,\%.
La serie incorpora además COVID-19 como categoría en 2021, por lo que el
cambio de composición debe reconocerse. No se interpreta como una medida
de ingresos, desempleo o victimización individual.

\begin{table}[htbp]
\centering\small
\caption{Indicadores de contexto económico y seguridad ciudadana}\label{tab:infosegura_contexto}
\begin{tabularx}{\textwidth}{@{}Yp{2.5cm}p{2.1cm}Y@{}}\toprule
Indicador & Período y unidad & Valor & Lectura para el estudio\\\midrule
Economía, principal problema & 2021; \% & 46,9 & Prioridad declarada, no ingreso observado\\
Seguridad, principal problema & 2021; \% & 26,6 & No es prevalencia de delitos\\
Privación por desempleo & 2019-2020; cambio en pp & +10,4 & Variación de un indicador de privación, no tasa de desempleo\\
Restricciones por inseguridad & 2019-2020; cambio en pp & $-6,8$ & Componente de privaciones; distinto de disciplina familiar\\
Inseguridad en transporte & 2021; \% reportado & 68,1 & Base por espacio no detallada en la fuente\\
Inseguridad en comunidad & 2021; \% reportado & 19,9 & Percepción, no hecho delictivo registrado\\\bottomrule
\end{tabularx}
\fuente{PNUD/Infosegura (2021). La base 18+ del indicador de uso de espacios no se atribuye automáticamente a los indicadores de percepción. Las unidades de las filas no se suman.}
\end{table}

El aumento del componente de privación por desempleo plantea preguntas sobre
recursos para la crianza, dependencia económica y acceso a servicios. Estas
relaciones son hipótesis; la variación de 10,4 puntos no demuestra que el
desempleo causara las prácticas observadas por MICS siete años antes.
La agenda de prevención de Infosegura también propone atender los espacios públicos de mayor
uso. Su aportación es ampliar la prevención hacia condiciones de acceso y
movilidad, junto con la protección dentro del hogar.

\subsubsection{Desplazamiento de prioridades y persistencia del temor}
Entre 2019 y 2021, la economía como principal problema aumenta 14,7 puntos,
mientras seguridad disminuye 33,2 puntos (figura \ref{fig:prioridades}). En
2021, COVID-19 representa 22,0\,\%. La incorporación de esa categoría cambia
el conjunto de respuestas: las variaciones de seguridad y economía no son
una redistribución exclusiva entre ambas. Tampoco miden la intensidad de
cada preocupación, pues elegir un problema principal puede coexistir con
preocupaciones importantes en otros ámbitos.

\begin{figure}[htbp]\centering
\begin{tikzpicture}\begin{axis}[width=0.85\textwidth,height=5.7cm,
 xmin=2016.8,xmax=2021.2,ymin=0,ymax=80,xtick={2017,2018,2019,2021},
 scaled x ticks=false,/pgf/number format/1000 sep={},ylabel={Porcentaje},
 grid=major,grid style={gray!15},tick label style={font=\small},
 legend style={at={(0.5,-0.18)},anchor=north,legend columns=3,column sep=6pt,draw=none,font=\small}]
\addplot[petroleo,thick,mark=*] table[x=anio,y=seguridad,col sep=comma]{data/infosegura_preocupaciones.csv};
\addlegendentry{Seguridad}
\addplot[ocre,thick,mark=square*] table[x=anio,y=economia,col sep=comma]{data/infosegura_preocupaciones.csv};
\addlegendentry{Economía}
\addplot[gris,only marks,mark=triangle*,mark size=3pt] coordinates {(2021,22.02)};
\addlegendentry{COVID-19}
\end{axis}\end{tikzpicture}
\caption{Evolución de las prioridades declaradas: principal problema del país}\label{fig:prioridades}
\fuente{PNUD/Infosegura (2021). Series seleccionadas, no distribución completa de respuestas. COVID-19 solo se observa como categoría en 2021; los años anteriores no se codifican como cero. Sin dato de 2020 ni IC propios.}
\end{figure}

La percepción de inseguridad también disminuye, pero conserva diferencias
marcadas por espacio (figura \ref{fig:percepcion}). Entre 2019 y 2021 cae
10,7 puntos en transporte público y 12,5 en comunidad. En 2021, las cifras
reportadas son 68,1\,\% y 19,9\,\%, respectivamente. Su distancia descriptiva
de 48,2 puntos muestra que un mismo diagnóstico nacional requiere observar
lugares de movilidad y de residencia. La fuente no detalla las bases de
cada espacio; la distancia no se interpreta como una comparación de riesgo
individual ajustada por exposición.

\begin{figure}[htbp]\centering
\begin{tikzpicture}\begin{axis}[width=0.85\textwidth,height=5.7cm,
 xmin=2016.8,xmax=2021.2,ymin=0,ymax=100,xtick={2017,2018,2019,2021},
 scaled x ticks=false,/pgf/number format/1000 sep={},ylabel={Porcentaje reportado},
 grid=major,grid style={gray!15},tick label style={font=\small},
 legend style={at={(0.5,-0.18)},anchor=north,legend columns=2,column sep=6pt,draw=none,font=\small}]
\addplot[petroleo,thick,mark=*] table[x=anio,y=transporte_publico,col sep=comma]{data/infosegura_percepcion.csv};
\addlegendentry{Transporte público}
\addplot[ocre,thick,mark=square*] table[x=anio,y=comunidad,col sep=comma]{data/infosegura_percepcion.csv};
\addlegendentry{Comunidad}
\end{axis}\end{tikzpicture}
\caption{Percepción de inseguridad en transporte y comunidad}\label{fig:percepcion}
\fuente{PNUD/Infosegura (2021). La base de personas de 18+ del indicador de uso de espacios no se traslada a estas series de percepción. Percepción distinta de victimización; segmentos entre años observados, sin imputar 2020.}
\end{figure}

La combinación de menor victimización declarada y persistencia del temor en
transporte sugiere evaluar varias dimensiones de seguridad a la vez. Para
las familias, el acceso a escuela, empleo y servicios requiere movilidad;
el temor puede ser relevante incluso sin una experiencia delictiva reciente.
Esta es una interpretación para orientar preguntas de investigación, no una
estimación del efecto del transporte sobre disciplina o primera unión. Se
necesitaría medir uso efectivo, frecuencia de desplazamiento y horarios para
examinar esa relación.

\subsubsection{Privaciones: cambios heterogéneos entre 2019 y 2020}
La figura \ref{fig:privaciones} examina los veinte cambios de privación disponibles.
Siete indicadores aumentan, once disminuyen y dos quedan sin cambio en la
precisión publicada. El aumento de privación por desempleo (+10,4 puntos)
destaca frente a otras variaciones positivas de hasta 1,3 puntos. Al mismo
tiempo, disminuyen restricciones por inseguridad ($-6{,}8$), falta de acceso
a salud ($-3{,}0$) e incidencia de crimen y delito ($-2{,}4$). El conjunto
documenta movimientos de distinto signo, relevantes para una interpretación
multidimensional del bienestar.

\begin{figure}[!htbp]\centering
\begin{tikzpicture}\begin{axis}[width=0.76\textwidth,height=11.0cm,xbar,
 bar width=7pt,xmin=-8,xmax=12,xtick={-8,-4,0,4,8,12},
 ymin=-0.7,ymax=19.7,ytick={0,...,19},
 yticklabels={Restricciones por inseguridad,Sin acceso a salud,Crimen y delito,
 Hacinamiento,Tenencia del terreno,Cuidado temprano inadecuado,
 Sin seguridad social,Sin saneamiento,Inasistencia escolar,
 Subempleo e inestabilidad,Trabajo infantil,Rezago educativo,
 Sin agua potable,Pared y piso inadecuados,Sin espacios de esparcimiento,
 Baja educación de adultos,Techo inadecuado,Inseguridad alimentaria,
 Daños y riesgos ambientales,Desempleo},
 xlabel={Variación 2019--2020 (puntos porcentuales)},
 axis x line*=bottom,axis y line*=left,xmajorgrids=true,grid style={gray!15},
 tick label style={font=\scriptsize},nodes near coords,
 nodes near coords style={font=\scriptsize,/pgf/number format/fixed,/pgf/number format/precision=1}]
\addplot[fill=petroleo,draw=none] table[x=cambio,y=pos,col sep=comma]{data/derivada_privaciones.csv};
\draw[gris,thick] (axis cs:0,-0.7)--(axis cs:0,19.7);
\end{axis}\end{tikzpicture}
\caption{Variación de los indicadores de privación}\label{fig:privaciones}
\fuente{PNUD/Infosegura (2021). Etiquetas abreviadas; textos originales en CSV. Cambios de indicadores de privación, no sus niveles. El componente desempleo no es la tasa de desempleo. Cero es ausencia de cambio al redondeo disponible.}
\end{figure}

La magnitud de desempleo invita a examinar apoyos de ingreso y cuidado junto
con servicios de protección. Sin embargo, no se suman las veinte variaciones
para obtener una variación de pobreza multidimensional: las privaciones
pueden coexistir en un hogar y el índice depende de ponderaciones, umbral y
distribución conjunta. Una reducción del indicador de inseguridad tampoco
implica que el nivel de privación sea bajo o que haya desaparecido. Sin
niveles iniciales, no pueden compararse niveles finales solo por la altura
de estas barras.

La lectura conjunta con MICS permite formular una agenda concreta: distinguir
estrés material, aceptación de prácticas violentas y exposición efectiva;
observar sus intersecciones con datos compatibles; y evaluar si los servicios
alcanzan a hogares con varias necesidades. Los agregados de 2019--2021 no
explican por sí solos la conducta familiar de 2014. Su contribución al
artículo consiste en precisar las dimensiones que una siguiente medición
debe integrar, manteniendo sus distintos períodos y poblaciones.

